% ---------------------------------------------------------------------------
% C.4'e EK — Prompt-framing stress testlerindeki ideoloji arka plan
% paragrafının tam metni. Appendix C.4'te Tablo 11-12'nin ALTINA yapıştır.
% §3.4'teki footnote bu etikete işaret etmeli: \ref{app:background}
% ---------------------------------------------------------------------------

\paragraph{Background paragraph used in the third stress test.}
\label{app:background}
The full text of the ideology background paragraph prepended to the prompt
in the third framing variant is reproduced below, with an English
translation.

\begin{quote}
Türkiye'nin ideolojik yapısı iki temel eksen üzerinde şekillenir: sol-sağ
siyasi ekseni ve seküler-dindar kültürel ekseni. Sol tarafta sosyalizm,
feminizm, sosyal demokrasi ve Kürt ulusal hareketi birbirine yakın
konumlanır; bu ideolojiler arasında geçişkenlik ve dayanışma yaygındır. Sağ
tarafta Türk milliyetçiliği ve kısmen de Kemalizm yer alır, ancak bu ikisi
kültürel eksende birbirinden ayrılır: Kemalizm güçlü biçimde sekülerken,
Türk milliyetçiliği dindarlığa daha yakındır. Muhafazakârlık ve İslamcılık
birbirine çok yakındır ve kültürel eksenin dindar ucunu oluşturur; ancak
İslamcılık siyasi eksende hafifçe sola kayar --- bu, özellikle Kürt
illerinde İslamcılığın sol hareketle tarihsel geçişkenliğini yansıtır. Kürt
ulusal hareketi sosyalizmle en yakın ideolojidir, ama aynı zamanda dindarlık
boyutunda ortalamanın üzerindedir; yani sol ama seküler değildir.
Türkiye'de üç temel ideolojik küme vardır: (1) merkez-sağ ve seküler küme
(Batı kıyıları, CHP tabanı), (2) merkez-sağ ve dindar küme (İç Anadolu,
Karadeniz, AKP-MHP tabanı), (3) sol ve dindar küme (Güneydoğu, Kürt nüfusun
yoğun olduğu bölgeler, DEM tabanı). Bu yapıda en büyük ideolojik yarık
Kemalizm-İslamcılık ve Türk milliyetçiliği-Kürt hareketi arasındadır.
\end{quote}

\emph{English translation.}
\begin{quote}
Turkey's ideological structure is shaped along two basic axes: a left-right
political axis and a secular-religious cultural axis. On the left,
socialism, feminism, social democracy, and the Kurdish national movement are
positioned close to one another; fluidity and solidarity among these
ideologies are common. On the right stand Turkish nationalism and, in part,
Kemalism, though the two diverge on the cultural axis: Kemalism is strongly
secular, whereas Turkish nationalism is closer to religiosity. Conservatism
and Islamism are very close to each other and form the religious end of the
cultural axis; Islamism, however, leans slightly left on the political axis
--- reflecting, especially in Kurdish provinces, its historical fluidity
with the left movement. The Kurdish national movement is ideologically
closest to socialism, but it is also above average on the religiosity
dimension; that is, left but not secular. There are three basic ideological
clusters in Turkey: (1) a center-right and secular cluster (the western
coasts, the CHP base), (2) a center-right and religious cluster (Central
Anatolia, the Black Sea, the AKP-MHP base), and (3) a left and religious
cluster (the Southeast, regions with dense Kurdish populations, the DEM
base). In this structure the deepest ideological rifts run between Kemalism
and Islamism and between Turkish nationalism and the Kurdish movement.
\end{quote}
